\subsection{Recovery Uncertainty and Apparent Source Information}
\label{app:recovery-information}

A local Gaussian model explains why preserving a recovery mean alone
can overstate confidence in a new source. Suppose, conditionally on an
earlier history,
\[
 r=Ox+sa+\epsilon,\qquad
 x\sim\mathcal N(0,P),\quad\epsilon\sim\mathcal N(0,R),
\]
where $x$ and $\epsilon$ are independent and $P,R$ are positive definite.
Here $R$ excludes the initial-state uncertainty represented by $P$.
Marginalizing that state gives $\Sigma(P)=R+OPO^\top$.
For a fixed local source direction $s$, Woodbury's identity gives
\begin{align}
 \mathcal J(P)&=s^\top\Sigma(P)^{-1}s\nonumber\\
 &=s^\top R^{-1}s
 -b^\top(P^{-1}+O^\top R^{-1}O)^{-1}b.
\end{align}
Here $b=O^\top R^{-1}s$. The nonnegative subtracted term measures how much of the apparent
source information can be explained by uncertain recovery coordinates.
If $P_2\succeq P_1$, then
$\mathcal J(P_2)\le\mathcal J(P_1)$ for this fixed direction and noise
law. This standard Gaussian identity explains a possible confidence
error; it does not prove that a larger covariance improves attribution.
Likelihood normalization, competing sources and false alarms still matter.

This decomposition is a separate explanatory model. The preceding channel
diagnostic already uses a covariance induced by the complete null
innovation model. Adding the bank's state covariance to it again would
count the same predictive uncertainty twice. Likewise, an uncertainty
already marginalized in $\Sigma$ must not also be fitted as a free
recovery direction without redefining the statistical model. The
implemented state--severity correlations are specified in
Section~\ref{app:conditional-moments}; independence is not assumed for
those stored posterior variables.

Port-Hamiltonian formulations provide a physical interpretation of
network interconnection and energy balance \cite{fiaz2013port}.
Their reduced swing-model energy is not automatically an energy function
for the present lossy DAE with exciters and governors. Similarly, the
input-change versus topology-change distinction is exact only within
its stated model: a finite impedance-load change can alter algebraic
admittance and the local linearization. The implemented load direction
is an affine approximation of that change. Energy residuals and
reconstructed hidden voltages are therefore not appended as independent
likelihood factors to the PMU reports that produced them.
